During the assessment, \team identified several key areas where \client demonstrated strong security practices. Continuation and further enhancement of these practices will significantly contribute to \clientP resilience against future threats, fostering long-term growth and maintaining user trust.

\begin{enumerate}
\item \textbf{Critical Risk Reduction in Safety \& Privacy Systems:} The most significant improvement since the last assessment is the successful remediation of high-liability vulnerabilities. \client has fixed the "Climate Control" system, which previously allowed for physical equipment takeover, and removed the "Galley" file share that exposed sensitive HR documents. By patching these critical areas, \client has directly reduced the risk of physical passenger harm and data privacy litigation, effectively protecting both the brand's reputation and its bottom line.\item \textbf{Hardened Defenses Around Ship Stability Operations:} The team has successfully reinforced the operating system supporting the Ballast Control System. Despite \teamP attempts to utilize stolen credentials or advanced technical attacks to breach the core server (SSH), the system successfully blocked unauthorized administrative access. This demonstrates a strong "Defense-in-Depth" strategy: even if outer layers are tested, the critical controls preventing catastrophic maritime accidents remain secure against administrative takeover.\item \textbf{Streamlined Infrastructure for Cost \& Risk Reduction:} The IT team has successfully consolidated the network environment by removing unnecessary servers, such as the standalone MySQL host, and migrating services to centralized locations. From a business perspective, this infrastructure hygiene significantly reduces the attack surface by limiting the number of potential entry points for attackers. This consolidation not only complicates adversary reconnaissance but also lowers the operational overhead associated with maintaining and securing redundant hardware.
\end{enumerate}
